\documentclass[11pt]{article}
\usepackage[a4paper,margin=2.5cm]{geometry}
\usepackage{amsmath,amssymb}
\usepackage{microtype}
\usepackage{enumitem}
\newcommand{\MNC}{M2NC}

\title{Report on ``Supersymmetric solutions of non-relativistic\\ 11-dimensional supergravity'' (C.~D.~A.~Blair)}
\date{}

\begin{document}
\maketitle

\section*{Summary of the paper}
Eleven-dimensional supergravity admits a non-relativistic ``membrane Newton--Cartan'' limit (\MNC{} SUGRA).
Consistency of supersymmetry in the limit requires constraints on the bosonic geometry, and earlier work found two consistent versions: a maximally supersymmetric one (32 parameters, with a heavy tower of constraints) and a ``half-maximal'' one (16 parameters, $\epsilon_-=0$, far fewer constraints).
The paper takes the \MNC{} limits of the M2 and M5 brane solutions, done in two ways.
In the \emph{scaled} limit the number of branes is rescaled by hand.
In the \emph{smeared} limit the brane is smeared over the longitudinal directions.
It then asks which version of the theory each limit solves and whether it is supersymmetric.
It also gives a simple longitudinal AdS$_3$ solution, which includes the BTZ black hole.

\section{Principal achievement}
The paper gives a clean, explicit and checked classification.
Scaled limits, which are localised in the longitudinal directions, solve the \emph{maximally} supersymmetric theory.
Smeared limits, which are localised in the transverse directions, have non-vanishing $f_{Aabc}$ and $T_a{}^{\{AB\}}$, so they solve only the \emph{half-maximal} theory.
All of these backgrounds are shown to be supersymmetric by solving the Killing spinor equations in closed form.
The solutions include the 8-spinor smeared M2 and M5 cases and an AdS-type 32-spinor solution for the near-horizon scaled M5.

The analysis also corrects and extends earlier results.
For the scaled M2, a general harmonic-in-$\ln H$ family of backgrounds is found with infinitely many Killing spinors, and 24 for generic $H$, where only 32 at the near-horizon point had been reported.
Several backgrounds show a non-Lorentzian enhancement of supersymmetry, in which arbitrary solutions of the 2d massless Dirac equation appear.
This is traced to the fermionic shift symmetries.
The resulting picture matches field-theory expectations for non-relativistic AdS/CFT: half of the supersymmetries become local, and there is an enhanced superconformal symmetry.
It also gives the otherwise odd half-maximal theory a concrete physical role.
It additionally suggests that the scaled/smeared distinction tracks maximal/half-maximal supersymmetry, which is a useful organising principle for the growing family of non-Lorentzian brane backgrounds.

\section{Principal drawback}
The central weakness is that everything is tested against a framework with no independent footing, so the conclusions are only as secure as that framework.
The constraints and supersymmetry rules of the \MNC{} theory are taken from earlier work that, as the author admits, relied partly on linearised evidence.
The half-maximal version in particular has a supersymmetry algebra that does \emph{not} close into diffeomorphisms.
Hence the claim that the smeared solutions are ``supersymmetric'' rests on the existence of Killing spinors alone.
That is a weak criterion in a setting where, as the paper itself says, there is no principled way to count preserved supercharges.
The ``infinite-dimensional'' enhancements are not fully counted.

A mild circularity follows.
The half-maximal theory is promoted as physically relevant because the smeared solutions fit it, but those solutions are the ones that fail the maximal constraints by construction.
The results are therefore consistency checks, not tests that could have falsified the theory.
Related points:
\begin{itemize}[nosep]
\item The longitudinal dependence of the scaled solutions is left unconstrained by the equations of motion, so they hold ``trivially''.
\item The by-hand scaling of $N$ stays unexplained, and no brane sources are included.
\item The Killing spinor analysis often rests on simplifying assumptions ($\partial_0\epsilon_-=0$, particular choices of $\eta_-$, $\rho_{+A}$), so completeness is not established.
\item The spherical M5 limit is acknowledged to be ``quick and dirty''.
\end{itemize}

\section{Most promising follow-up research idea}
Balancing feasibility and impact, I would pursue \textbf{coupling \MNC{} SUGRA to the non-relativistic brane worldvolume action}, that is, including source terms, starting in the simplest tractable setting.
The paper itself flags this in its summary.

\paragraph{Concrete programme.}
Reduce the \MNC{} limit of the M2 along a worldvolume direction to the D2NC (Matrix 2-brane theory) string setting.
There the string worldsheet action in this limit exists in the literature, and the supergravity limit descends by dimensional reduction from 11d.
Then:
\begin{enumerate}[nosep]
\item Check the coupled equations $S_{\text{SUGRA}}+N S_{\text{worldvolume}}$ directly for the smeared and scaled backgrounds. Does a source term reproduce the smeared harmonic function? Is an extra scaling of $N$ ever needed or consistent?
\item Compare the supersymmetry of the probe, via $\kappa$-symmetry, with the Killing spinors found here. This would give an independent test of the claimed half-maximal structure.
\item If the picture holds, extend to the M5 and, via reduction and duality, to the D-brane limits used in non-relativistic AdS/CFT proposals.
\end{enumerate}

\paragraph{Why this one.}
It is feasible because the bosonic check is a finite computation in a setting where all ingredients (action, reduced supergravity, explicit solutions) already exist.
It is also high-impact, because it attacks the paper's main weakness directly.
It would supply the first independent, source-based criterion for which solutions belong to which version of the theory.
It could explain the scaled-versus-smeared puzzle, since the scaling of $N$ may prove inconsistent with a source while smearing is natural.
It would also test whether the constraint tower is physically right.
Cheaper but less decisive alternatives are checking closure of the supersymmetry algebra on the Killing spinors found here, and the Freund--Rubin-type extension suggested in a footnote.
Full classification of Killing spinors would be more ambitious and rests on the same shaky framework.

\end{document}
